\section*{ANNEXE A – patati patata}
\addcontentsline{toc}{section}{Annexe A – patati patata}

Votre rapport peut compter un nombre indéterminé d’annexes, numérotées A, B, C, qui ne sont pas décomptées dans le nombre de pages du rapport. Étant optionnelles, les évaluateurs peuvent ne pas les lire. C’est l’occasion de placer des documents qui rompraient la continuité du texte principal, tels que des résultats d’expériences de vérification, ou qui l’allongeraient trop. Par exemple, plutôt que de mettre trois graphes pour présenter les résultats sur trois échantillons différents, ce qui peut prendre 1.5 pages, on n’en conserve qu’un dans le texte principal pour présenter brièvement la méthode, ainsi qu’un tableau de synthèse, les courbes pour les deux autres échantillons étant renvoyées en annexe.



\section*{Annexe : Planning du stage par phases hebdomadaires}
\addcontentsline{toc}{section}{Annexe B - Planning de l'alternance}

Le stage s’est déroulé du \textbf{12 mai au 25 juillet 2025}, soit une durée totale de 11 semaines. Le tableau suivant présente la répartition hebdomadaire des différentes phases du stage:

\bigskip


\begin{ganttchart}[
    hgrid,
    vgrid,
    x unit=0.9cm,
    y unit title=0.7cm,
    y unit chart=0.9cm,
    bar height=0.7,
    bar/.style={fill=blue!30, rounded corners=3pt},
    milestone/.style={fill=orange!50},
    bar label font=\footnotesize\bfseries,
    milestone label font=\footnotesize\bfseries,
    group label font=\small\bfseries,
    title label font=\small\bfseries,
    group/.style={draw=black, fill=green!20, rounded corners=3pt},
    group height=0.5,
    group peaks={0}{0}{0.3},
    bar left shift=0.1,
    bar right shift=-0.1
]{1}{11}
\gantttitle{Semaines du stage}{11} \\
\gantttitlelist{1,...,11}{1} \\

\ganttgroup{Phase préparatoire}{1}{3} \\
\ganttbar{Accueil et cadrage}{1}{1} \\
\ganttbar{Configuration environnement}{1}{1} \\
\ganttbar{Étude de l'existant}{2}{3} \\
\ganttmilestone{Milestone 1: Fin d'étude}{3} \\

\ganttgroup{Phase de tests}{4}{9} \\
\ganttbar{Tests d'intrusion}{4}{6} \\
\ganttbar{Analyse des résultats de l'audit}{7}{9} \\
\ganttmilestone{Milestone 2: Tests complets}{8} \\

\ganttgroup{Phase finale}{9}{11} \\
\ganttbar{Correction des vulnérabilités}{9}{10} \\
\ganttbar{Rédaction rapport}{10}{11} \\
\ganttmilestone{Milestone 3: Fin de stage}{11}
\end{ganttchart}
